% Isolated post-CP33 candidate. No canonical or shared-input mutation.
\subsection{Compatible cyclic intermediates on one cubic slice}
\label{subsec:compatible-cyclic-intermediates}

Work over an algebraically closed field of characteristic two. Let
\((B,\Theta_B)\) be a nonordinary principally polarized abelian
surface, put \(N=\mathcal O_B(3\Theta_B)\), and let
\[
 \pi:A\longrightarrow B,\qquad
 D=\ker\pi\simeq(\mathbf Z/\ell\mathbf Z)^2,
 \qquad L=\pi^*N,
 \tag{CCI1}
\]
where \(\ell\) is an odd prime. The isogeny is etale. Let
\(U\subset\mathbf P(H^0(A,L))\) be the smooth, sharp open,
\(\mathcal V\) its relative \(\mathbf F_2\)-Lang local system,
and \(r=9\ell^2-1\) its rank. The full-radical theorem gives
the canonical proper sub-local system \(\mathcal K\subset\mathcal V\)
corresponding to the unique maximal proper geometric submodule
\(K\). The statements below use the arbitrary odd-source
Cartier lemma and quotient bound with their full-image hypotheses
exactly as stated.

For each line \(H\subset D\), set
\[
 A_H=A/H,\qquad q_H:A\longrightarrow A_H,
 \qquad\eta_H:A_H\longrightarrow B,
 \qquad M_H=\eta_H^*N.
 \tag{CCI2}
\]
Then \(\pi=\eta_Hq_H\), both degrees are \(\ell\),
\(q_H^*M_H=L\), and \(h^0(A_H,M_H)=9\ell\).
These are exact pullbacks of line bundles, not merely numerical
polarization equalities.

\begin{lemma}[The line-averaging identity]
\label{lem:compatible-line-averaging}
\status{Candidate; complete direct proof supplied for independent audit.}
For a subgroup \(J\subset D\), write
\(e_J=\sum_{g\in J}g\in\mathbf F_2[D]\).
Then
\[
 e_J^2=e_J,\qquad
 \sum_{\substack{H\subset D\\ |H|=\ell}}e_H=1+e_D,
 \qquad e_De_H=e_D.
 \tag{CCI3}
\]
Consequently, in every \(\mathbf F_2[D]\)-module \(V\),
\[
                  \bigcap_{|H|=\ell}\ker(e_H:V\to V)=0.
 \tag{CCI4}
\]
No splitting field for the characters of \(D\) is required.
\end{lemma}
\begin{proof}
Every element of \(J\) occurs \(|J|\) times in the product
\(e_J^2\), and \(|J|\) is odd. Every nonidentity element of
\(D\) belongs to exactly one line, whereas the identity belongs
to \(\ell+1\) lines, an even number. This proves the first
two identities. Each element of \(D\) occurs \(\ell\)
times in \(e_De_H\), proving the third. If every \(e_H\)
kills \(v\), their sum gives \(v=e_Dv\), and any one line
gives \(e_Dv=e_De_Hv=0\). Thus \(v=0\).
\end{proof}

\begin{proposition}[A common full-rank slice for every intermediate]
\label{prop:compatible-intermediate-common-slice}
\status{Candidate; complete direct proof supplied for independent audit.}
There is a nonempty open
\(S_D\subset\mathbf P(H^0(B,N))\) whose pullback lies in \(U\)
and, for every line \(H\), in the full smooth, sharp source slice
\(\mathbf P(H^0(A_H,M_H))\). On this one open the deck action
of \(D\) acts on the actual local system \(\mathcal V|_{S_D}\).
The operator \(e_H\) is the restriction of the averaging
projector on the \(H\)-invariant source slice; its image is the
relative Lang local system of the intermediate curve on \(A_H\).

Suppose, for a particular \(H\), that the source comparison
proves \(\mathcal K\subseteq\ker e_H\) on a nonempty open of
that intermediate source slice. Then the same inclusion holds on
\(S_D\). This assertion does not require the first open to meet
\(S_D\), and does not compare full fibres through a rank drop.
\end{proposition}
\begin{proof}
Decompose \(\pi_*\mathcal O_A\) into its odd character lines
on \(B\). Choose a smooth, sharp degree-one section of \(N\)
where all nontrivial character multiplication blocks are invertible.
The common Cartier incidence theorem makes this a nonempty open
condition in the degree-one coefficient space. Denote that open by
\(S_D\). Its pullback divisor on \(A\) is smooth because
\(\pi\) is etale, and connected by ample connectedness.
Its relative Lang rank is \(r\), so the resulting section belongs
to \(U\).

The characters of \(\eta_H\) are precisely the characters of
\(D\) trivial on \(H\). Thus their nontrivial multiplication
blocks are among the blocks already made invertible. The section on
\(A_H\) is likewise smooth, connected and sharp. The remaining
\(\pi\)-character blocks form the cosets belonging to the
nontrivial \(q_H\)-characters. Their invertibility is exactly
the invertibility of the outer blocks after the intermediate
decomposition; taking odd halves merely permutes the character
cycles. This proves the claimed simultaneous full-rank conditions.

For clarity, let \(T_H\) be the open in
\(\mathbf P(H^0(A_H,M_H))\) on which the source is smooth and
sharp, the pulled-back section belongs to \(U\), and all outer
character cycles are bijective. It is a nonempty open in a
projective space, hence integral and connected, and contains
\(\eta_H^*S_D\). The odd-source twisted-block lemma shows
directly that the extra determinant conditions are genuine open
conditions; the preceding character comparison supplies the stated
common points.

Over \(T_H\), deck translations by \(H\) preserve the
universal pulled-back divisor and commute with Frobenius. Their
averaging projector therefore acts on the actual
\(\mathbf F_2\)-local system. Odd etale descent identifies its
image with the source Lang system, including the relative quotient:
odd pullback identifies the constant ambient cohomology. On
\(S_D\), this action is literally the restriction of the
\(D\)-deck action, since \(H\subset D\) and all the divisors
are pullbacks of the same divisor on \(B\).

Finally, the composition
\[
 \mathcal K|_{T_H}\ \longrightarrow\
 \mathcal V|_{T_H}\ \xrightarrow{\ e_H\ }\
 \mathcal V|_{T_H}
 \tag{CCI5}
\]
is a morphism of finite etale local systems. If it vanishes on a
nonempty open, it vanishes on connected \(T_H\): after a finite
etale trivialization it is a locally constant matrix, and its
zero locus is open and closed. Restriction to \(S_D\) now gives
the asserted inclusion. There is no choice of separate generic-fibre
identifications for different \(H\): all these morphisms are
restrictions of the single sub-local system \(\mathcal K\)
on \(U\).
\end{proof}

\begin{theorem}[Full intermediate images force the full target image]
\label{thm:compatible-cyclic-intermediate-criterion}
\status{Candidate; conditional on the stated intermediate full images.}
Assume that, for every line \(H\subset D\), the relative
geometric monodromy of the complete system \(|M_H|\) on \(A_H\)
is
\[
                 \operatorname{GL}_{9\ell-1}(\mathbf F_2).
 \tag{CCI6}
\]
Then \(K=0\), and the relative geometric monodromy of the
complete system \(|L|\) on \(A\) is
\[
                 \operatorname{GL}_{9\ell^2-1}(\mathbf F_2).
 \tag{CCI7}
\]
\end{theorem}
\begin{proof}
For each \(H\), apply
Proposition~\ref{prop:DE-general-odd-source-quotient-bound}
to the inner cover \(\eta_H:A_H\to B\), of degree
\(a=\ell\), and the outer cover \(q_H:A\to A_H\), also
of degree \(\ell\). Its explicit full-source hypothesis is
exactly (CCI6). It gives \(\mathcal K\subseteq\ker e_H\)
on a nonempty open in the intermediate coefficient space.
The preceding proposition transports each such inclusion to the
same open \(S_D\). At one geometric point of that open, all
the operators are the operators in the same group algebra
\(\mathbf F_2[D]\) on the same fibre of \(\mathcal V\).
Equation (CCI4) therefore gives \(\mathcal K|_{S_D}=0\).
Its rank is constant on connected \(U\), so \(K=0\).
The full-radical image theorem, specifically
Corollary~\ref{cor:root-full-radical-submodules-orbits}, now
gives (CCI7).
\end{proof}

\begin{corollary}[A prime-degree criterion with type \((1,\ell)\) sources]
\label{cor:compatible-cyclic-principal-prime-criterion}
\status{Candidate; no intermediate full-image hypothesis suppressed.}
Let \((A,\Theta_A)\) be a fixed nonordinary principally polarized
abelian surface, and choose an outgoing maximal-isotropic
\(\ell\)-quotient as in (CCI1), with exact principal descent
\(\pi^*\mathcal O_B(\Theta_B)=\mathcal O_A(\ell\Theta_A)\).
For each of its \(\ell+1\) lines, let
\(T_H=\eta_H^*\mathcal O_B(\Theta_B)\).
The polarization \(T_H\) has type \((1,\ell)\) and
\(M_H=T_H^{\otimes3}\). If the complete cubic systems of
these particular \(\ell+1\) polarizations have the full
relative images (CCI6), then the fixed target system
\(|3\ell\Theta_A|\) has full relative image.
In particular a uniform full-image theorem for every odd cyclic
pullback of a cubic principal system of degree \(\ell\) would
imply the fixed-principal-system full-image theorem in degree
\(\ell\).
\end{corollary}
\begin{proof}
The exact outgoing descent is the one already used for the odd
principal systems. The line \(T_H\) is ample and
\(T_H^2=2\ell\). Its polarization type \((d_1,d_2)\)
satisfies \(d_1d_2=\ell\) and \(d_1\mid d_2\); since
\(\ell\) is prime, the type is \((1,\ell)\). Its cube is
\(M_H\), and the composite pullback is exactly
\(\mathcal O_A(3\ell\Theta_A)\). Apply the theorem.
\end{proof}

\begin{remark}[The remaining cyclic-source gate]
\label{rem:compatible-cyclic-intermediate-scope}
\status{Candidate scope; no new unconditional degree asserted.}
This criterion uses compatible intermediate quotients of one fixed
deck group. Distinct maximal-isotropic deck groups and their
potentially disjoint eigenvector slices are not compared. The
common-point argument takes place entirely in the smooth, sharp
locus and uses locally constant morphisms; it requires no
specialization isomorphism through a rank drop.
The full images in (CCI6) remain unproved in the new nonprincipal
cases. For a single cyclic cover, the known central slice gives
only \(K\subseteq\ker e_D\). Its nontrivial character
multiplicities and the dimensions of irreducible
\(\mathbf F_2[D]\)-modules alone do not force \(K=0\).
No numerical experiment or new external theorem is used here.
\end{remark}
